\documentclass[11pt]{article}
\usepackage[margin=1in]{geometry}
\usepackage{amsmath, amssymb, amsthm}
\usepackage{hyperref}

\title{The Math Behind Market Making:\\From Glosten-Milgrom to Avellaneda-Stoikov}
\author{Abhi Wadhwa}
\date{}

\begin{document}
\maketitle

\begin{abstract}
These are my notes on the mathematics of market making, covering the two models that I think give the clearest picture of what's going on. I start with Glosten-Milgrom (1985), which explains \emph{why} bid-ask spreads exist (adverse selection), and then move to Avellaneda-Stoikov (2008), which tells you \emph{how} to optimally set your quotes when you care about inventory risk. The connection between the two is that the first model answers ``why is there a spread at all?'' and the second answers ``how wide should my spread be, and where should I center it?''
\end{abstract}

\section{Why Do Spreads Exist?}

Before getting into the math, let me explain the core intuition. A market maker stands ready to buy and sell an asset. They quote a bid (buy price) and an ask (sell price), and pocket the difference on round-trip trades. Sounds like free money, right?

The catch is \textbf{adverse selection}. Some of the people you trade with know more than you do. When an informed trader buys from you, the true value is probably above your ask---you just sold too cheaply. When they sell to you, the true value is probably below your bid---you just bought too expensively. You systematically lose to these people.

The spread exists to compensate for these losses. You earn the spread on trades with uninformed (noise) traders, and you lose on trades with informed traders. In equilibrium, these balance out. \textbf{The key insight is that the spread is the ``price'' of adverse selection.}

\section{Glosten-Milgrom (1985)}

Here's the simplest version of the model, and honestly it's kind of elegant.

\subsection{Setup}

An asset has an unknown true value $V$, which is either $V_H$ (high) or $V_L$ (low). The prior probability that $V = V_H$ is $\mu$ (start with $\mu = 1/2$). There are two types of traders:
\begin{itemize}
\item \textbf{Informed traders} (fraction $\alpha$): they know $V$ and trade accordingly---buy when $V = V_H$, sell when $V = V_L$.
\item \textbf{Noise traders} (fraction $1 - \alpha$): trade randomly, buying or selling with equal probability regardless of $V$.
\end{itemize}

The market maker doesn't know who they're facing on any given trade. They have to set bid and ask prices to break even in expectation.

\subsection{Computing the Ask}

The ask is the price at which the market maker is willing to sell. If someone comes in to buy, the market maker asks: what's the expected value of the asset \emph{given that I'm seeing a buy order}?

Using Bayes' rule:
\[
P(V = V_H \mid \text{buy}) = \frac{P(\text{buy} \mid V = V_H) \cdot P(V = V_H)}{P(\text{buy})}
\]

Now let's compute each piece. If $V = V_H$, then informed traders buy (probability $\alpha$) and noise traders buy with probability $\frac{1-\alpha}{2}$. So:
\[
P(\text{buy} \mid V = V_H) = \alpha + \frac{1-\alpha}{2} = \frac{1+\alpha}{2}
\]

If $V = V_L$, informed traders sell (not buy), so only noise traders buy:
\[
P(\text{buy} \mid V = V_L) = \frac{1-\alpha}{2}
\]

With $\mu = 1/2$:
\[
P(\text{buy}) = \frac{1}{2} \cdot \frac{1+\alpha}{2} + \frac{1}{2} \cdot \frac{1-\alpha}{2} = \frac{1}{2}
\]

So the posterior is:
\[
P(V = V_H \mid \text{buy}) = \frac{\frac{1+\alpha}{2} \cdot \frac{1}{2}}{\frac{1}{2}} = \frac{1+\alpha}{2}
\]

And the competitive ask price is:
\[
\text{Ask} = E[V \mid \text{buy}] = \frac{1+\alpha}{2} \cdot V_H + \frac{1-\alpha}{2} \cdot V_L
\]

\subsection{Computing the Bid}

By symmetric reasoning:
\[
P(V = V_H \mid \text{sell}) = \frac{1-\alpha}{2}
\]
\[
\text{Bid} = E[V \mid \text{sell}] = \frac{1-\alpha}{2} \cdot V_H + \frac{1+\alpha}{2} \cdot V_L
\]

\subsection{The Spread}

The bid-ask spread is:
\[
\text{Spread} = \text{Ask} - \text{Bid} = \alpha \cdot (V_H - V_L)
\]

This is a beautiful result. \textbf{The spread is directly proportional to the fraction of informed traders} ($\alpha$) \textbf{and the size of the uncertainty} ($V_H - V_L$). If there were no informed traders ($\alpha = 0$), the spread would be zero. If everyone were informed ($\alpha = 1$), the spread would equal the full range of possible values, and the market maker would earn zero expected profit regardless.

Notice that at these competitive prices, the market maker earns zero expected profit. They make money on noise traders (who buy above and sell below the midpoint) and lose it to informed traders (who only trade when it's favorable). The spread exactly balances these.

\subsection{Bayesian Updating Over Time}

In practice, the market maker updates their belief $\mu$ after each trade. After observing a buy:
\[
\mu' = P(V = V_H \mid \text{buy}) = \frac{1+\alpha}{2} \cdot \frac{\mu}{P(\text{buy})}
\]

After a sell:
\[
\mu' = P(V = V_H \mid \text{sell}) = \frac{1-\alpha}{2} \cdot \frac{\mu}{P(\text{sell})}
\]

So the belief drifts toward the truth over time---buy orders push $\mu$ up, sell orders push it down. The spread narrows as $\mu$ approaches 0 or 1 (less uncertainty remaining). This is price discovery in action.

\section{Avellaneda-Stoikov (2008): Optimal Market Making}

Glosten-Milgrom tells us why spreads exist, but it doesn't tell us how to \emph{optimally} manage quotes when we care about inventory risk. For that, I turn to Avellaneda and Stoikov.

\subsection{Setup}

The midprice follows arithmetic Brownian motion:
\[
dS = \sigma \, dW
\]
(no drift under the risk-neutral measure). The market maker holds inventory $q$ and has a CARA utility function with risk aversion $\gamma$:
\[
U(x) = -e^{-\gamma x}
\]

The market maker quotes a bid $S - \delta^b$ and ask $S + \delta^a$, where $\delta^b$ and $\delta^a$ are the half-spreads on each side. Orders arrive as Poisson processes with intensity that depends on the spread:
\[
\lambda(\delta) = A \cdot e^{-k\delta}
\]
where $A$ is the baseline arrival rate and $k$ controls how sensitive traders are to the spread (wider spread $\Rightarrow$ fewer orders).

\subsection{The Reservation Price}

Here's where it gets interesting. The trick is that a market maker with inventory doesn't value the asset at the midprice $S$. If I'm holding a long position and the price drops, I lose money. The more inventory I hold, the more I want to shade my quotes to reduce that inventory.

Avellaneda and Stoikov show that the \textbf{reservation price}---the price at which the market maker is indifferent to a small trade---is:
\[
\boxed{r = S - q \cdot \gamma \cdot \sigma^2 \cdot (T - t)}
\]

This is one of those results that's easy to see once you understand it but took me a while to fully appreciate. If $q > 0$ (long inventory), $r < S$---the market maker values the asset below mid, because they're exposed to downside risk. If $q < 0$ (short), $r > S$. The adjustment scales with risk aversion $\gamma$, volatility $\sigma^2$, and time remaining $T - t$ (more time = more risk).

\subsection{Optimal Spread}

The optimal total spread (distance between bid and ask) is:
\[
\boxed{\delta^* = \gamma \sigma^2 (T-t) + \frac{2}{\gamma} \ln\!\left(1 + \frac{\gamma}{k}\right)}
\]

The first term is the \textbf{inventory risk premium}---it increases with volatility and risk aversion. The second term captures the \textbf{market impact / adverse selection component}---it depends on how elastic order flow is (parameter $k$).

\subsection{Putting It Together: Skewing Quotes}

The optimal bid and ask are:
\begin{align*}
\text{Bid} &= r - \frac{\delta^*}{2} = S - q\gamma\sigma^2(T-t) - \frac{\delta^*}{2} \\
\text{Ask} &= r + \frac{\delta^*}{2} = S - q\gamma\sigma^2(T-t) + \frac{\delta^*}{2}
\end{align*}

The midpoint of the quotes is the reservation price $r$, not the market midprice $S$. This means the quotes are \textbf{skewed} based on inventory:
\begin{itemize}
\item Long inventory ($q > 0$): both quotes shift down, making it more likely you sell (reducing inventory).
\item Short inventory ($q < 0$): both quotes shift up, making it more likely you buy.
\end{itemize}

This is exactly what real market makers do. After playing with this for a while in the simulator, I found that even a simple linear skew (shift quotes by $-q \times c$ for some constant $c$) dramatically improves PnL stability compared to symmetric quotes.

\section{Connecting the Two Models}

Glosten-Milgrom and Avellaneda-Stoikov are complementary:

\begin{center}
\begin{tabular}{lll}
& \textbf{Glosten-Milgrom} & \textbf{Avellaneda-Stoikov} \\
\hline
Focus & Why spreads exist & How to set optimal spreads \\
Key risk & Adverse selection & Inventory risk \\
Spread driver & Informed trader fraction $\alpha$ & Volatility $\sigma$, risk aversion $\gamma$ \\
Inventory & Not modeled & Central to the solution \\
Result & Zero-profit equilibrium & Utility-maximizing quotes \\
\end{tabular}
\end{center}

In reality, both forces operate simultaneously. A market maker faces adverse selection \emph{and} inventory risk. The Glosten-Milgrom spread is a lower bound on the spread you'd want to charge (to avoid being picked off by informed traders), and the Avellaneda-Stoikov framework adds the inventory management layer on top.

\begin{thebibliography}{9}
\bibitem{glosten1985} Glosten, L.R., Milgrom, P.R. (1985). ``Bid, Ask and Transaction Prices in a Specialist Market with Heterogeneously Informed Traders.'' \textit{Journal of Financial Economics}, 14(1), 71--100.

\bibitem{kyle1985} Kyle, A.S. (1985). ``Continuous Auctions and Insider Trading.'' \textit{Econometrica}, 53(6), 1315--1335.

\bibitem{avellaneda2008} Avellaneda, M., Stoikov, S. (2008). ``High-Frequency Trading in a Limit Order Book.'' \textit{Quantitative Finance}, 8(3), 217--224.

\bibitem{ohara1995} O'Hara, M. (1995). \textit{Market Microstructure Theory}. Blackwell Publishers.
\end{thebibliography}

\end{document}
